\documentclass[a4paper, 10pt]{article}

\usepackage[swedish]{babel}
\usepackage[T1]{fontenc}
\usepackage[utf8]{inputenc}
\usepackage{hyperref}
\usepackage{csquotes}
\usepackage[backend=biber, style=apa, sorting=nyt]{biblatex}
\usepackage{pdfpages}
%\usepackage{geometry}
\usepackage{xcolor}
\usepackage{soul}
\usepackage{fancyhdr}



\addbibresource{bibliografi.bib}

\newcommand{\mytitle}{HP601A---Examinationsuppgift}
\newcommand{\myauthor}{Calle Ketola}

\author{\myauthor}
\title{\mytitle}

\setcounter{tocdepth}{2}

%\setlength{\headheight}{22.5pt}

\begin{document}

\fancyhead[c]{\mytitle}
\fancyhead[r]{\myauthor}
\fancyhead[l]{Akademiskt lärande}

\pagestyle{fancy}

\pagenumbering{roman}

\maketitle

\tableofcontents

\clearpage

\pagenumbering{arabic}

\section{Inledning}

I detta paper diskuteras ett kursmoment i
en kurs som nyligen getts vid Malmö 
universitet. Momenten analyseras utifrån 
hur det relaterar till tre punkter, 
konstruktiv länkning, likabehandling samt 
styrdokument samt tre perspektiv, student-,
lärar- och samhällsperspektiv. Det ges 
även ett förslag på hur detta moment kan 
designas om för att bättre nå lärandemålen 
för kursen och utbildningen som helhet.

Det kursmoment som jag har valt att analysera 
är den första workshopen i kursen 
\textit{CD102A Objektorienterad 
programmering} som läses i läsperiod två 
det första läsåret utav programmet 
\textit{Civilingenjör i Datateknik} vid Malmö 
universitet. Anledningen till att jag valt 
detta moment är för att när vi genomförde 
momentet hösten 2025 blev studenterna 
stressade och workshopen drog över tiden med 
ungefär 45 minuter. En anledning till att 
workshopen drog över tiden var för att 
bedömningarna tog längre tid än planerat. 

En anledning till att jag velat förändra 
uppgiften är att om man följer Feisel-Schmittz 
taxonomi så ligger ''att bara lösa uppgiften'' 
lågt på skalan av enkelhet och komplexitet 
\parencite[sida 159]{elmgrenolding2025}. 
Målet i kursen, och utbildningen, är att 
studenterna ska kunna analysera och utvärdera 
olika lösningar \parencites{cd102a2025}{tcdat2025}
vilket ligger överst i Feisel-Schmittz taxonomi
vilket är vad jag vill jobba mot i detta moment.


\section{Det valda kursmomentet}

I det valda momentet fick studenterna i uppgift 
att två-och-två skapa ett program utifrån en 
uppgiftsbeskrivning, samt skapa ett klassdiagram
som modellerade deras program. I programmet 
skulle de skapa tre klasser, klassen 
\texttt{Galaxy}, klassen \texttt{Star},
samt klassen \texttt{Planet}. Till sin hjälp att 
göra detta fick studenterna beskrivningar av 
varje klass, vilka attribut de skulle innehålla, 
samt korta beskrivningar över vilka operationer 
klasserna skulle ha.

För att kontrollera att de var färdiga hade 
studenterna fått en \textit{driver}-kod som 
testade funktionaliteten i deras lösningar. 
När den skrev ut ''\textit{Present your 
solution}'' så fick de lov att redovisa sin 
lösning för en lärare.

Under workshopen blev ca 90~\% av de deltagande
studenterna godkända. Då bedömningarna drog 
över tiden, är det oklart hur många som 
faktiskt var färdiga när passets egentliga 
stopptid hade nåtts. För att bli godkända skulle 
studenterna kunna förklara sina lösningar och 
visa upp ett korrekt klassdiagram som visade 
hur klasserna relaterade till varandra.

Under ett uppföljande moment veckan efter fick 
studenterna möjlighet att bygga vidare på sina 
lösningar från detta moment. Då var det flera 
som hade problem att göra ändringar i sin kod 
vilket tyder på att det studenterna hade gjort 
under det tidigare momentet inte hade befästs 
tillräckligt för att jobba vidare med \parencite[sida 22]{elmgrenolding2025}. 
Detta ser jag som ett bevis för att studenternas 
kunskap här ligger lågt i Feisel-Schmittz 
taxonomi \parencite[sida 159]{elmgrenolding2025}.

\subsection{Konstruktiv länkning}

%Hur har detta relaterat till senare moment i 
%kursen? Särskilt inlämningar och tentamen? Hur 
%relarar det till tidigare moment?

Detta är kursens första bedömda moment, och 
sätter således tonen för kommande moment. 
Inför detta moment har studenterna haft fem 
föreläsningar och fyra övningstillfällen. 
De första två övningstillfällena har behandlat 
skillnader mellan programspråken Java, som de 
ska använda i den här kursen, och Python, som 
de har använt i den föregående kursen 
\textit{CD100A Imperativ programmering}. 
Under tredje tillfället fick studenterna en 
uppgift där klassbeskrivningarna har getts på 
ett liknande sätt som i momentet vi ska 
undersöka, och under det fjärde 
övningstillfället har de fått använda 
ett verktyg för att skapa klassdiagram.

Den här workshopen bygger således vidare på 
övningstillfälle tre och fyra. Skillnaderna är
dels att de nu ska göra en kombination av de 
båda tidigare övningarna samt att programmet 
testar sig själv. För att kontrollera att 
programmet var redo att presenteras skulle 
studenterna köra den \textit{driver}-kod som 
de fick till uppgiften. Denna kod utförde 
tester på studenternas kod och skrev ut vad 
det försökte göra, det förväntade svaret, 
samt det erhållna svaret från studentens 
kod. Detta lät studenterna snabbt se om något 
inte var rätt, på ett ungefär var det var fel, 
och ungefär vad som var fel. Detta möjliggjorde 
och förenklade till viss del en självbedömning 
hos studenterna \parencite[sida 278]{elmgrenolding2025}.
Det underlättade även för lärarna då 
studenterna i samma utsträckning inte behövde 
be läraren kontrollera om svaret var korrekt. 
Rättningen i slutet av passet underlättades också 
eftersom programmet berättade att det fungerade 
som det skulle.

I de två kommande bedömda workshoparna är den 
tematiska inramningen den samma, med stjärnor 
och planeter, men uppgifterna ändrar karaktär.
De andra workshoparna saknade självrättande 
delar och mängden kod som behövde skrivas var 
mindre. Att mängden kod som behövde skrivas 
minskades berodde framförallt på att den första 
workshopen drog över på tiden. Således var den 
konstruktiva länkningen inte särdeles stark 
till dessa moment.

Kopplingen var desto starkare till de två 
bedömda inlämningsuppgifterna som följde. 
Framförallt eftersom de tekniker de uppmuntrades 
att använda i det första momentet kunde 
användas för att lösa problem som uppstod i de 
båda inlämningsuppgifterna. Det var särskilt 
tekniker för att hämta information ur objekt 
som låg i andra objekt, samt för att hantera 
arrayer med objekt som de hade användning av.

När tentamen kom så återanvände vi delar av 
modellen för stjärnor och planeter som vi 
använt i workshopen. Detta ledde till att flera 
studenter gav mer utförliga svar än vad som 
krävdes för godkända poäng.

Ett problem som kan ha uppstått är att eftersom 
det här momentet blev så stressigt är att 
studenterna leddes in i en mentalitet där 
själva skrivandet av kod är det viktigaste och 
att kvantitet --- som i att skriva kod snabbt ---
är viktigare än kvalitet --- som i att skriva 
rätt kod och utvärdera valet av lösning.

\subsubsection{Studentperspektiv}

För studenter som har som mål att bli 
civilingenjörer och få ett ''bra'' jobb, 
de som av John Biggs symboliseras av Robert
\parencite[sida 16]{biggstang2022}, har 
uppgiften uppskattats. Detta vet vi efter 
samtal med studenter efter kursens slut.
För studenterna som är mer som Susan, och 
mer är ute efter kunskapen för kunskapens 
skull, har uppgiften varit stressande på 
grund av att de inte har hunnit förstå 
hur deras lösningar fungerar. Sen har vi en
tredje kategori studenter som består av 
dem 10 \% som inte blev godkända. För dessa 
studenterna var uppgiften för stor och 
för svår. Det är oklart om de hade 
klarat uppgiften om de fått ännu mer tid 
än de extra 45 minuterna som det blev. Att 
några studenter inte klarar uppgiften är 
inte fel i sig, särskilt som inte alla i 
gruppen var godkända på alla moment från 
föregående kurser.

\subsubsection{Lärarperspektiv}

När uppgiften skapades var tanken att den 
skulle:

\begin{enumerate}
   \item knyta samman det som gjorts i kursen hittills,
   \item knyta an till den senaste föreläsningen,
   \item vara lätt att bedöma.
\end{enumerate}

Uppgiften var således inte nämnvärt framåtriktad, 
istället var fokus på vad som hade gjorts 
tidigare. Detta kan ge en konstruktiv länkning 
till tidigare material, men inte nödvändigtvis 
till lärandemålen. Dock knyter uppgiften an till 
lärande mål sju, samt till viss del till 
lärandemål fyra och åtta \parencite{cd102a2025}.
På så vis kan man se en konstruktiv länkning 
mellan uppgiften och lärandemålen.

Uppgiften blev inte svår att bedöma, men den tog 
tid att bedöma. Att avgöra om en student förstår 
sin lösning är en tidskrävande process. Samtidigt
kan studentens förståelse fördjupas när hen 
får sätta ord på vad hen har gjort 
\parencite[sida 175]{elmgrenolding2025}.

\subsubsection{Samhällsperspektiv}

Under den här kursens gång har en kritik mot 
konstruktiv länkning som lyfts varit att om 
innehållet i föreläsningar och övningar för 
starkt kopplas till examinerande moment, 
alltså om man undervisar det som ska examineras 
och examinerar det så som det har undervisats,
så finns det en risk att studenterna får en 
smal kunskap och en falsk tilltro till 
instruktioner.

En mem inom lärarkåren är studenter och elever 
som ställer frågan ''kommer det här på 
provet/tentan''. Skulle läraren svara nej så 
tappar studenterna intresset för det som 
diskuteras. I stunden kan det för studenten vara 
ett rationellt beslut att inte ta emot den 
nya informationen för att minska mängden 
information som behöver bearbetas. Men i längden 
går studenten miste om information som hade 
kunnat vara viktig utanför examinationens ramar. 
Detta är ett tecken på att studenterna ser 
undervisningen som ett verktyg för att få betyg,
betyget är sen ett verktyg för att få examen,
och examen ett verktyg för att få ett jobb. 
Detta indikerar att studenterna ser 
högskolan som en \textit{kunskapsfabrik} där 
skolan ska producera så många jobbredo individer 
som möjligt \parencite[sida 274]{branden2022}.
En möjlig risk med detta är att högskolan 
producerar individer som kan lösa problem som 
är snarlika det dem har stött på under sin 
utbildning, men inte problem som är ''nya''.

Att lägga ifrån sig pennan när studenten får 
höra att det inte kommer på tentan är också ett 
tecken på att studenten ägnar sig åt 
ytinriktat lärande 
\parencite[sida 22]{elmgrenolding2025}. Detta 
är inte målet med utbildningen, istället vill 
vi har studenter som istället för att ha 
memorerat fakta och lösningsmetoder har 
förstått hur saker hänger ihop och hur det 
fungerar, så kallat djupinlärning 
\parencite[sida 22]{elmgrenolding2025}.
 
\subsection{Likabehandling}

Som civilingenjörsstudenter har studenterna med 
sig kurserna Fysik 1 och Fysik 2 från gymnasiet,
eller Komvux. I dessa kurser får de bland annat 
lära sig om universum och astronomi 
\parencite{kursplanfysik2011}. Det borde därför 
inte vara något i bakgrunden till uppgiften som 
inte alla redan har en tidigare förståelse för.
Studenterna har också läst en introducerande 
kurs i programmering som de förväntas ha goda 
kunskaper från, även om de ännu inte är 
godkända. 

Det är fler kvinnor än män som söker till 
högskoleutbildningar \parencite{uhrsokande2025},
men utbildningar med inriktning på ''data'' har 
färre kvinnliga sökande än män (se bilaga 
\ref{app:antagna}). Det är ungefär lika stor 
andel av de antagna männen som kvinnorna 
som klarar en utbildning med inritning på 
elektronik, datateknik, eller automation 
(se bilaga \ref{app:examen}). Jag ser inte heller 
något i hur uppgiften läggs upp som skulle 
gynna eller missgynna antingen män eller kvinnor. 
Uppgiften är också frånkopplad från frågor rörande 
människor i allmänhet. Den mänskliga kopplingen 
i uppgiften är att studenterna förväntades att 
jobba i par, där de själva fick välja vem de 
skulle jobba med. Skulle några studenter bli 
över matchade vi lärare ihop dessa med varandra. 
Denna matchning skedde genom att vi parade ihop 
de studenter som stod närmast varandra. Vi 
beaktade således inte studenternas kön, 
etniska tillhörighet, eventuella 
funktionsvariationer eller andra saker som är 
diskrimineringsgrundande.

\subsubsection{Studentperspektiv}

Studenter vill som regel inte särbehandlas. 
Detta gäller både positiv och negativ 
särbehandling. Det här innebär bland annat att 
du inte vill få mer eller mindre hjälp än dina 
medstudenter. Detta kan vara problematiskt, 
särskilt om studenten behöver mer hjälp än sina 
kamrater. Samtidigt har somliga studenter rätt 
till särskilt stöd, så som tolk eller inläst 
material.

Ett sätt att läsa detta på är genom att 
genomföra sin undervisning på ett stöttande sätt. 
På gymnasiet är det vanligt att elever med 
neuropsykologiska funktionsvariationer får extra 
tydliga instruktioner eller checklistor som 
de kan bocka av. Detta är något man som lärare 
kan ge alla studenter, eftersom detta är något 
som de allra flesta blir hjälpta av.

För mycket av den sortens hjälp skulle däremot 
kunna hindra studentens progression. För att 
inte det ska ske kan man bygga in det i sin 
undervisningsplan. Då kan man i början av 
undervisningen delge färdiga checklistor, för 
att senare tillsammans med studenterna skapa 
checklistor, så att de till sist kan skapa 
checklistorna själva.

\subsubsection{Lärarperspektiv}

Som lärare finns det en svårighet i att använda 
sina fördomar och förutfattande meningar om en 
grupp på ett konstruktivt sätt. Det är också 
svårt att göra skillnad på en fördom och en 
generaliserad verklighet. Ett exempel på detta 
är att om jag ser en student som använder en 
mac så tror jag inte att hen förstår sig på 
filträd i datorn. Skulle studenten istället 
använda Linux så skulle jag förutsätta att 
studenten hade bra koll på filträd. Båda fallen 
är forutfattade meningar som jag har om 
studenter, där i alla fall den första har 
motbevisats.

Men den här typen av fördomar mot studenterna 
kan göra att man lägger undervisningen på fel 
nivå. Risken är stor att man lägger sig på en 
för lätt nivå när man pratar med en 
mac-användare eller en för hög nivå när man 
pratar med en Linux-användare. För att minimera 
problemet behöver man som lärare vara lyhörd och 
snabbt kunna växla mellan nivåer, och då helst 
på ett sätt som inte riskerar att skämma ut 
studenten eller dig själv.

\subsubsection{Samhällsperspektiv}

Att alla har lika stor rätt till utbildning 
slås fast av de mänskliga rättigheterna 
\parencite[sida 7]{unhumanrights}. Under kursens
gång har detta problematiserats inom 
utbidlningar som leder till legitimation, 
exempelvis sjuksköterske- och 
receptiarieutbildningarna. Detta är utblidningar 
som är tänkta att leda till jobb där den 
anställde förväntas kunna utföra ett arbete, 
och där funktionsnedsättningar kan göra 
detta arbete omöjligt att utföra i praktiken.
En blind sjuksköterska skulle ha svårt 
att göra bedömningar på exempelvis hudsjukdomar.

Om vi ser på högskolan som en kunskapsfabrik så 
är detta problematisk 
\parencite[sida 274]{branden2022}. 
Här har vi en student med  rätten att studera 
men som inte kommer att kunna bli godkänd på 
flera av sina kurser, och därför inte kan nå 
examen. Ser vi istället högskolan som ett 
\textit{samhällsförpliktigande universitet} så 
är den blinda studentens omöjlighet att bli 
legitimerad sjuksköterska inte ett problem. Studenten 
kan fortfrande ta till sig kunskap och växa 
som individ och samhällsmedborgare 
\parencite[sida 283]{branden2022}. Eventuellt 
skulle även en blind sjuksköterskestudent kunna 
bidra med insikter som andra studenter och 
lärare inte hade tänkt på annars.

\subsection{Policy och styrdokument}

De policy och styrdokument som används i det 
här papret är framförallt kursplanen samt 
utbildningsplanen.

\subsubsection{Kopplingar till kursplan}

I kursplanen står det att laborationerna ska 
bedöma målen:

\begin{enumerate}
    \item förklara grundläggande begrepp inom objektorienterad programmering
    \item redogöra för skillnader mellan imperativ (procedurell) och objektorienterad programmering, samt för- och nackdelar med olika programspråksparadigm,
    \item förklara och ge exempel på olika sätt att organisera samlingar av objekt
    \item konstruera mindre objektorienterade program där polymorfism nyttjas för att ge objekt generaliserat beteende \parencite{cd102a2025}.
\end{enumerate}

I det här momentet så testas framförallt mål 1 
och det förbereder för att mål 4 ska kunna visas 
senare. Anledningen till att det inte riktigt kan 
visa på färdigheter inom mål 4 är för att 
generalisering och polymorfism inte hunnits tas 
upp i kursen än. 

\subsubsection{Kopplingar till utbildningsplan}

Ett av målen i utbildmningplanen är ''visa 
förmåga att planera och med adekvata metoder 
genomföra kvalificerade uppgifter inom givna 
ramar'' \parencite{tcdat2025}. Vilket studenterna 
jobbar mot i det här momentet.

\subsubsection{Studentperspektiv}

Momentet kan kopplas till både kursplanen och 
utbildningsplanen, detta har dock inte tydligt 
kommunicerats till studenterna. Istället har 
de fått information om att detta är ett 
obligatoriskt moment. Detta är oturligt, då 
utmärkt undervisning har tydliga mål för såväl 
lärare som studenter \parencite[sida 142]{branden2022}
Detta tvingar studenterna att blint lita på 
sina lärare och tro på att de har planerat 
momentet utifrån lärandemålen.

\subsubsection{Lärarperspektiv}

I skapandet av uppgiften har det som tidigare 
skrivits inte varit lärandemålen som direkt har 
format uppgiften. Istället har uppgiften formats 
efter tidigare moment. Nu har det råkat sig att 
uppgiften ändå knyter an till lärandemålen fyra,
sju och åtta vilka ska examineras under det här 
momentet.

För lärarna som ska bedöma och handleda under 
momentet är det bra att veta vad 
uppgiftsförfattaren vill lära ut eller testa 
studenterna på \parencite[sida 142]{branden2022}.
Det gör det lättare att handleda och styra 
studenterna i rätt riktning.

\subsubsection{Samhällsperspektiv}

Ur ett samhällsperspektiv är det viktigt att 
man har en uppfattning om vad studenterna kan 
när de är klara med en kurs eller utbildning.
Detta är bland annat för att göra det lättare 
för studenter att flytta mellan universitet 
och för att lättare kunna anställas av 
näringslivet \parencite{uhrbologna}. För att 
detta ska vara fallet behöver 
examinationsmomenten bedöma det som står i 
kursplanen.

\section{Det nya kursmomentet}

Det nya kursmomentet går ut på att 
studenterna får tre olika koder som ger 
samma resultat när de körs. Det tre 
programmen är skrivna i programspråket 
Java och hanterar stjärnssystem. När 
programmen körs kan användaren visa 
information om två förinlagda stjärnssystem,
Sol samt Proxima Centauri, med deras 
omkretsande planeter. Användaren kan 
också lägga till nya stjärnsystem och 
fylla dessa med nya planeter som 
programmet sedan kan presentera.

%% Behöver jag gå så här djupt i varje 
%% del?

Den första lösningen är programmerad med 
en imperativ lösning, alltså helt utan att 
använda objektorienterade lösningar. All 
information sparas i arrayer med strängar 
och kopplingen mellan stjärnor och planeter 
sker genom att en stjärna och dess planeter 
ligger på samma index i olika arrayer.

Den andra lösningen är en generell lösning 
med objektorientering. Det finns en klass 
för själva programmet och användargränssnittet, 
samt en klass som representerar en stjärna, 
och en klass som representerar en planet. 
Användargränssnittet har en array med stjärnor,
och stjärnorna har varsin array med planeter. 
För att presentera information om stjärnorna 
och planeterna hämtar gränssnittet informationen 
direkt från objekten.

Den tredje lösningen följer 
\textit{Model-View-Controller} designprincipen
där användergränssnittet inte har någon direkt 
koppling till himlakropparna. Istället 
förmedlas informationen via en mellan-klass som 
hanterar programmets logik. Här har mellanklassen
en array med stjärnor --- som i sin tur har arrayer
med planeter --- och som instruerar 
användargränssnittet att presentera viss data.

Studenterna ska läsa igenom de tre programkoderna
och köra programmen. De ska sedan diskutera 
skillnader och likheter i de tre lösningarna och 
presentera vilken lösning de tycker är bäst.
Förhoppningen är att studenterna kommer fram till 
att det inte finns någon lösning som är objektivt
bättre än de andra utan att alla tre har sina 
olika för- och nackdelar.

För att bli godkända så ska studenterna delta 
aktivt under hela passet, vilket kontrolleras 
löpande av närvarande lärare, samt kunna redogöra 
för olika för- och nackdelar mellan de olika 
lösningarna.

\subsection{Förändringar}

I detta avsnittet tas förändringarna av uppgiften 
upp samt hur dessa gör att uppgiften bättre 
kopplas till lärandemålen för både kursen och 
programmet. 

\subsubsection{Ändrad arbetsform}

Uppgiften har gått från att studenterna jobbar
två och två med att skriva kod, till att de 
sitter i grupper om tre och diskuterar kod.
Det är tre olika anledningar till att denna 
förändring har gjorts.

Den första anledningen är att de skriver kod 
under de övriga laborationspassen. Under dessa 
passen kan de be om hjälp utan att det finns en 
press om att uppgiften ska bedömas. Även om 
det var tillåtet att be om hjälp i den tidigare 
uppgiften var det få som vågade göra det 
eftersom det var en examinerande uppgift. Vi 
kontrollerar även deras förmåga att skriva kod
med andra examinationsuppgifter, så att kolla 
det nu är överflödigt. Det var dessutom 
hyfsat enkelt att be en ''smart kompis'' på 
internet, generativ AI, att generera en lösning.
Ber man Copilot365 att jämföra koderna och ge 
för- och nackdelar med de tre olika programmen 
får man tyvärr ett bra svar även med den nya 
uppgiften.

Den andra anledningen är rättningsbördan. Vi 
upptäckte att redovisningarna tog väldigt lång 
tid när vi genomförde det ursprungliga momentet.
Detta trots att programmet var självrättande i 
den mån att studenterna fick ett meddelande om 
koden gjorde det vi förväntade oss. Att gå igenom 
deras lösningar och se så att de även förstod 
vad de hade gjort tog så pass lång tid att passet 
drog över med ungefär en 45 minuter. Att istället 
kunna gå runt och lyssna på studenterna under 
passet för att sedan höra deras slutsatser i 
slutet kan gå betydligt snabbare.

Den tredje och största anledningen till att 
byta arbetsform är att jag anser att studenterna 
behöver få se bra kod. De behöver få se korrekta 
lösningar med en kodstandard vi förväntar oss att 
de följer. Detta ger studenterna en tydlig 
förebild vilket kan gynna deras fortsatta 
inlärning \parencite[91]{elmgrenolding2025}  .
Får studenterna inte se vad vi menar är en bra 
lösning är det svårt för dem att själva avgöra 
kvalitén i sina egna och kamraters lösningar 
vilken kan minska värdet av framtida 
kamratgranskning \parencite[sida 280]{elmgrenolding2025}. Vi vill även att de börjar 
fundera på sina val av lösningar. Genom att 
diskutera tre olika alternativ är förhoppningen 
att studenterna får en medvetenhet om att det 
finns olika sätt att lösa ett problem och att 
de framgent gör aktiva val av metoder för att 
lösa sina egna problem. 

\subsubsection{Förtydliganden i uppgiften}

En förändring i hur uppgiften presenteras 
är att det nu finns med ett stycke om den
pedagogisk idéen bakom momentet. Att låta 
studenterna veta att det finns en medveten 
plan kring momentets upplägg stärker 
uppgiften på två sätt.

Först försäkras studenten om att det 
finns en plan och att det inte är ett 
godtyckligt moment som finns med ''bara för 
att''. Det kan ge en tillit till 
undervisningen och lugna studenten med att 
om hen följer lärarens plan och instruktioner 
så kommer hen att komma vidare i sitt lärande
\parencite[sida 123]{branden2022}.

Den andra sättet detta stärker uppgiften på 
är att det visar studenten hur momentet 
kopplas till lärandemålen och varför vi som 
lärare anser att detta är en bra metod för att 
nå dit. Detta är särskilt viktigt när man 
avviker från vad många ser som det traditionella 
sättet att lära sig kodning, nämligen 
\textit{learning by doing}. Det är synnerligen 
viktigt då jag själv ofta repeterar för 
studenterna att knacka kod är nödvändigt för 
att bli bättre på att koda. Att nu förklara 
varför uppgiften avviker från ett etablerat 
mönster kan minska motståndet från studentera 
att ta sig an uppgiften ordentligt 
\parencite[sida 160]{branden2022}.

\subsection{Konstruktiv länkning}

Att jämföra olika program är något som kan 
testas på en tentamen. Där kan studenterna få två 
eller fler olika koder som ska jämföras där 
de sedan ska diskutera för och nackdelar med 
de olika implementationerna. I inläming två 
ställs också studenterna inför designbeslut och 
förhoppningen är att detta moment har förberett 
dem för att fatta det beslutet.

Faller utfallet ut väl i det här momentet skulle 
följande workshopar i kursen kunna följa samma stuk, 
där de får färdig kod som ska analyseras och 
behandlas på någotvis. Det skulle även vara 
möjligt att inlämningsuppgifterna ska redovisas 
på liknande sätt där studenterna får läsa igenom 
varandras lösningar och diskutera sina olika 
designval. Annars, som det ser ut nu, är där 
inte någon nämnvärd konstruktiv länkning till 
övriga moment i kursen.

Om man inte gör andra ändringar i övriga 
kursmoment hamnar vi i en sits där uppgiften i 
sig är en bättre uppgift, men den konstruktiva 
länkningen minskar.

\subsection{Likabehandling}

I den tidigare arbetsformen har studenter med 
dålig datorvana, dyslexi, eller andra hinder 
som gör att de har svårt att skriva på ett 
tangentbord missgynnats. När text ska skapas 
på en begränsad tid kan det orsaka en stress 
och få studenter att underprestera jämfört 
med om uppgiften varit utformad på ett annat 
sätt.

Den nya varianten av uppgiften missgynnar 
fortfarande dyslektiker eftersom det blir mer 
att läsa. Att få kod uppläst är problematisk 
av framförallt två anledningar:

\begin{enumerate}
    \item Det är \textit{konstiga} ord och ingen 
    naturlig meningsbyggnad. Ofta förkortas ord 
    något godtyckligt och språk kan blandas. 
    Exempelvis är uttrycket 
    ''\texttt{System.out.println(arr[i]/5 == 3);}''
    ett fullt logiskt och syntaktiskt korrekt 
    uttryck i Java. Men för en textuppläsare är 
    det svårt att läsa upp korrekt. Detta kan 
    exempelvis testas i Microsoft Word som har 
    en textuppläsningsfunktion. Är Word inställt 
    på svenska läser den upp ''System out println
    arr i snedstreck fem lika med lika med tre''.
    Medan en programmerare skulle kunna läsa det 
    som ''Skriv ut: är arr på plats i, delat på fem 
    lika med tre?''.
    \item Textuppläsning tar tid. De tar som 
    regel 30~\% längre tid att läsa en text högt 
    än vad det tar för en genomsnittlig person 
    att läsa en text tyst \parencite{BRYSBAERT2019104047}. 
    Men med den rapporterade minskade 
    läsförståelsen kanske detta är ett glapp 
    som numera minskat \parencite{sydsvenskan241111}.
\end{enumerate}

Den tidigare uppgiften var --- inklusive kod 
som behövde läsas --- tio sidor lång. Den nya 
uppgiften är --- med koden inräknad --- ca 30 
sidor lång. Detta är en markant ökning av 
text. Däremot har producerandet av text strukits. 
Detta kan kompensera för den större textmassan 
som ska läsas. Även om det totala antalet sidor 
som ska läsas har ökats har uppgiftsbeskrivningen 
dragits ner från sex sidor till två och en halv 
sidor.

Att läsa kod kan inte riktigt likställas med 
att läsa en löpande text. En programkod som 
huvudsakligen är skriven på engelska kan 
snarare likställas med strukturerad engelska
(\textit{eng. Structured English}). En sida 
programkod innehåller ca 200 ord (detta är 
taget från koden studenterna får i den här 
uppgiften). Varje rad på en sida är alltså 
förhållandevis kort vilket underlättar för 
personer med lässvårigheter 
\parencite{dyslexiascotland2025}. Jämförelsevis
har det här papret cirka 500 ord per sida.

Den nya utformningen kan försvåra för studenter 
som inte har svenska som modersmål \parencite{balter2023}. Ett krav 
för att komma in på utbildningen 2025 är ett 
godkänt betyg i Svenska 3 eller Svenska som 
andraspråk 3 \parencite{tcdat2025}. Att föra en 
ämnesteknisk diskussion är svårt om man har låga 
ämneskunskaper --- vilket de flesta studenter har 
i början av sin utbildning --- det är ännu 
svårare om man också har svaga språkkunskaper.
Detta är något som studenterna behöver öva på 
vilket de får möjlighet till i den här uppgiften.

\subsection{Policy och styrdokument}

Liksom tidigare version av momentet så undersöks 
den nya versionen jämtemot kursplanen och 
utbildningsplanen.

\subsubsection{Kopplingar till kursplan}

De lärandemål i kursplanen som momentet är menat 
att, framförallt, träna är 

\begin{itemize}
    \item Efter avslutad kurs skall studenten kunna: 
    redogöra för skillnader mellan 
    imperativ (procedurell) och objektorienterad 
    programmering, samt för- och nackdelar med 
    olika programspråksparadigm.
    \item Efter avslutad kurs skall studenten kunna:
    motivera val av lämpliga kontrollstrukturer 
    och klasser vid konstruktion av ett 
    objektorienterat program.
\end{itemize}

Ett problem här är att som kursplanen är 
skriven idag så ska laborationerna inte bedöma 
det första utav lärandemålen. Det här passet 
förbereder däremot studenterna för att kunna 
visa den kunskapen när de kommer till den 
skriftliga tentamen som ska bedöma det 
lärandemålet \parencite{cd102a2025}.

Uppgiften skapar också en brygga mellan den 
tidigare kursen CD100 och den här kursen genom 
att visa en lösning som följer principerna från 
den kursen, och visar en lösning som följer 
principerna i den här kursen. Den kan även ses 
som förberedande för att i andra moment nå andra 
kursmål.

\subsubsection{Kopplingar till utbildningsplan}

Ett av målen för civilingenjörsutbildningen 
på Malmö universitet är att studenten ska 
utveckla sin ''förmåga att göra självständiga 
och kritiska bedömningar'' \parencite{tcdat2025}. 
Den nya utformningen av uppgiften riktar in 
sig på det här målet genom att här låta 
studenterna få kritiskt jämföra olika 
kodstycken. 

Enligt högskoleförordningen ska ''studenten visa 
kunskap om det valda teknikområdets vetenskapliga 
grund och beprövade erfarenhet samt insikt i 
aktuellt forsknings- och utvecklingsarbete''
\parencite{hogskoleforordningen2025}. I uppgiftens 
nya utformning får studenten ta del av tre olika 
paradigmer inom programmering som är en del av 
proffessionens gemensamma erfarenhet.

Vad samhället förväntar sig av en färdigutbildad 
civilingenjör framgår i högskoleförordningen.
Där framgår bland annat att en civilingenjör ska 
kunna ta hänsyn till samhälleliga aspekter, samt 
miljö- och arbetsmiljöaspeketer i sitt arbete 
\parencite{hogskoleforordningen2025}. 
Att kunna värdera sina lösningar mot vad en 
situation kräver leder till kod som är lättare 
att underhålla och vidareutveckla. Det leder 
till mjukvara som fungerar en större del av tiden 
och som lättare kan lagas när det behövs.

Att känna till olika lösningar och paradigm 
gör också att det blir lättare att arbeta med 
olika projekt, olika projektgrupper, och mot 
olika mål. På ett sätt kan man kalla en 
utvecklare med en bred kunskap av olika 
paradigm för ''mångkulturerad'', i den mening 
att hen är kunnig och delaktig inom flera 
olika programmeringskulturer. 

\section{Hur AI har använts i den här texten}

I skrivprocessen för den har rapporten har 
AI använts för att formatera .bib-filen med 
referenser korrekt. Exempelvis har Copilot365 
använts för att få fram vilka fält som ska 
fyllas i när man refererar till lagstiftning.
Copilot365 har även använts för att kontrollera om 
texten är sammanhängande, om det är några 
brister i det framförda resonemanget, samt om 
språkbruket är konsekvent. 

% Följande promptar har använts:

% \begin{itemize}
%     \item ''Jag håller på och skriver ett paper 
%     och behöver referera till lagar och 
%     förordningar. Hur lägger jag in 
%     Högskoleförordningen i en bib-fil?''
%     \item ''Tack. Hur lägger jag in en kursplan 
%     publicerad av ett universitet i min bib-fil?''
%     \item ''Enligt Karolinska institutet ska 
%     kursplaner skrivas ut så här i 
%     referenslistan när man använder APA: 
%     "Universitet. (År). Kursplan. Titel. 
%     Universitet, Institution. URL"
%  När jag typsätter min fil får jag följande 
%  format: "Titel. (år). Universitet. URL"

%  I min tex-fil har jag angett: \begin{verbatim}\usepackage[backend=biber, style=apa, sorting=nyt]{biblatex}\end{verbatim} och min bib-fil ser ut som följer:
% \begin{verbatim}@online{
%     cd102a2025,
%     organization = {Malmö universitet},
%     institution = {Teknik och samhälle},
%     url = {https://utbildningsinfo.mau.se/kurs/kursplan/
%         sv/a3294619-8619-11ef-9561-318924f24352/20252},
%     title = {Kursplan för CD102A: Objektorienterad 
%         programmering},
%     year = {2025},
%     urldate = {2026--02--05}
% }”\end{verbatim}
%     \item ''Hur påverkar fältet 'langid' referensen?''
%     \item ''Hur skrivs en tidningsartikel in i en bib-fil?''
% \end{itemize}

%Jag skulle vilja ha hjälp att analysera texten ur följande synvinklar:
%1. Är texten sammanhängande?
%2. Är det några hål i resonemangen som förs?
%3. Är språkbruket konsekvent?

\printbibliography[heading=bibintoc]

\clearpage




\appendix

\section{Antagna och sökande hösten 2025}
\label{app:antagna}

\begin{center}
\begin{tabular}{|l|rrr|}    
    \multicolumn{4}{c}{Sökande och antagna efter andra urvalet hösten 2025}\\
    \hline
                    & Män   & Kvinnor   & Totalt\\ \hline
    Antal sökande   & 73149 & 33487     & 106636\\
    Andel sökande   & 69 \% & 31 \%     & 100 \%\\ \hline
    Antal antagna   & 5663  & 2182      & 7845\\
    Andel antagna   & 72 \% & 28 \%     & 100 \% \\ \hline
\end{tabular}

Källa: \parencite{uhrantagna2026}
\end{center}

Sökningen som gjordes var program som innehåller 
''data'' som nyckelord. Detta är en bred 
sökning som kan få med utbildningar som inte 
har programmering. Dessutom skulle utbildningar 
som saknar en ''data''-tagg men som ändå 
innehåller programmering saknas. Men detta 
bortfall tordes vara försumbart, likaså den 
extra datan.

\section{Genomströmning på program med inriktning Elektronik, datateknik och automation}
\label{app:examen}

\begin{center}
\begin{tabular}{|l|rr|}
    \hline
            & Män       & Kvinnor \\ \hline
    Antagna & 10171     & 2040\\
    Examen  & 3556      & 694\\
    Andel   & 35 \%     & 34 \%\\ \hline
\end{tabular}
\end{center}

Tabellen visar hur många män respektive kvinnor 
tagit examen från ett Civil- eller 
högskoleingenjörsprogram med en inriktning mot 
elektronik, datateknik, eller automation som 
inom tre år efter att deras nominella studietid
har gått. Datan visar studenter som påbörjat sina 
ingenjörsstudier mellan åren 2012 och 2018 
\parencite{scbexamen2026}.



\clearpage

\section{Tidigare version av uppgiften}

\includepdf[pages=1-]{Bilagor/CD102A-ht25-workshop-1.pdf}

\section{Uppdaterad version av uppgiften}

\includepdf[pages=1-]{Bilagor/CD102A-HT26-Worskhop-1.pdf}

\end{document}